% SHARD-ID: AQM-43-FINITE-FIELD-SCALE-CONTINUUM-LIMITS
% SHARD-TITLE: Finite-Field Scales and Continuum Limits
% SHARD-SUMMARY: Defines base-extension scales \(X_r=X_0\times_{\mathbb F_q}\mathbb F_{q^r}\) and their closed-point Weyl nets.
% SHARD-SUMMARY: Proves the trace obstruction to naive Weyl embeddings and the trace-normalized replacement.
% SHARD-SUMMARY: Computes the point \(X_0=\operatorname{Spec}\mathbb F_3\), its prime-to-\(3\) continuum limit, and Frobenius action.
% SHARD-KEYWORDS: continuum limit, finite field extension, base change, Weyl net, trace, Frobenius

\section{Finite-Field Scales and Continuum Limits}
\label{sec:finite-field-scale-continuum-limits}

\subsection{Status and Source Anchors}

This shard adds finite-field extension rank as the first arithmetic
``scale'' parameter.  The closed-point Weyl net is AQM-28 and the dynamics
language is AQM-42.  The scheme facts use the Stacks Project snapshot:
\path{references/algebraic_geometry/stacks_project_algebra.tex}, lines
250--274, 2972--2983, 3104--3117, 3310--3350, 6520--6665, and 45278--45284.
Frobenius is sourced from
\path{references/algebraic_geometry/stacks_project_varieties.tex}, lines
7005--7112.  The finite-field trace formula also appears in
\path{references/heisenberg_weil/sidana_kashyap_2202_00248_source/EAQECCs_over_rings_v15.tex},
lines 789--801.

\subsection{Lamport Dependency Skeleton}

\[
\begin{array}{rcl}
D1&:&k=\mathbb F_q,\ k_r=\mathbb F_{q^r}\subset\bar k,\
      k_r\subset k_s\Longleftrightarrow r\mid s.\\
D2&:&X_r=X_0\times_{\Spec k}\Spec k_r,\
      \pi_{s,r}:X_s\to X_r.\\
D3&:&\mathcal A_r(U)\text{ is the AQM-28 closed-point Weyl net on }X_r.\\
L1&:&\text{A closed point }x\in X_r\text{ splits over }k_s\text{ through }
      B_x=\kappa(x)\otimes_{k_r}k_s.\\
L2&:&\text{The naive diagonal label map multiplies the trace commutator by }
      n=s/r.\\
L3&:&\text{A trace-one element }c_x\in B_x\text{ gives a Weyl-compatible
      label embedding.}\\
T1&:&\text{Trace-normalized local embeddings induce monomorphisms of
      open-local Weyl nets.}\\
T2&:&\text{Prime-to-}p\text{ scales have canonical coherent embeddings.}\\
T3&:&X_0=\Spec\mathbb F_3\text{ gives }
      \varinjlim_{3\nmid r}M_{3^r}(\mathbb C).\\
T4&:&\text{Frobenius acts by }W(q,p)\mapsto W(q^3,p^3).
\end{array}
\]

\subsection{Scales and Closed-Point Nets D1--L1}

Fix an algebraic closure \(\bar k\) of \(k=\mathbb F_q\), and let
\(k_r=\mathbb F_{q^r}\) be the unique subfield of \(\bar k\) with \(q^r\)
elements.  Then \(k_r\subset k_s\) exactly when \(r\mid s\).  For an affine
\(k\)-scheme \(X_0=\Spec R\), define the rank-\(r\) scale
\[
  X_r=X_0\times_{\Spec k}\Spec k_r=\Spec(R\otimes_k k_r).
\]
If \(r\mid s\), base change along \(k_r\subset k_s\) gives
\(\pi_{s,r}:X_s\to X_r\).

For an open \(U\subset X_r\), \(\mathcal A_r(U)\) denotes the AQM-28
closed-point Weyl algebra:
\[
  \mathcal A_r(U)=
  \varinjlim_{S\Subset U\cap |X_r|_{\rm cl}}
  \bigotimes_{x\in S}\operatorname{End}_{\mathbb C}\ell^2(\kappa(x)).
\]
Thus a scale change should compare \(\mathcal A_r(U)\) with
\(\mathcal A_s(\pi_{s,r}^{-1}U)\).

Let \(x\in |X_r|_{\rm cl}\), put \(K=\kappa(x)\), and write
\(B_x=K\otimes_{k_r}k_s\).
Since finite fields are separable,
\(B_x\cong\prod_{y\mapsto x}\kappa(y)\), where \(y\) ranges over the closed
points of \(X_s\) above \(x\).  This is the
finite algebraic form of site splitting at finer scale.

\subsection{The Trace Obstruction L2}

The one-site character at \(x\) is
\[
  \psi_x(a)=\exp\!\left(\frac{2\pi i}{p}\Tr_{K/\mathbb F_p}(a)\right).
\]
The naive diagonal map sends \((q,p_0)\in K^2\) to
\(\Delta_x(q,p_0)=((q_y,p_{0,y}))_{y\mapsto x}\),
where \(q_y,p_{0,y}\) are the images in \(\kappa(y)\).  Let
\(a=p_0q'-q p'_0\).
The fine-scale commutator phase of the diagonal images is
\[
  \prod_{y\mapsto x}
  \exp\!\left(\frac{2\pi i}{p}
  \Tr_{\kappa(y)/\mathbb F_p}(a_y)\right).
\]
The exponent is
\[
  \Tr_{B_x/\mathbb F_p}(a\otimes1)
  =
  \Tr_{K/\mathbb F_p}\!\left(\Tr_{B_x/K}(a\otimes1)\right).
\]
As a \(K\)-algebra, \(B_x\) has dimension \(n=s/r\), and \(a\otimes1\) acts
on it as scalar multiplication by \(a\).  Hence
\[
  \Tr_{B_x/K}(a\otimes1)=n a,
\]
so the fine phase is the coarse phase raised to the \(n\)-th power:
\[
  \prod_{y\mapsto x}\psi_y(a_y)=\psi_x(a)^n .
\]
Therefore the naive diagonal label map is Weyl-compatible for all labels only
when \(n\equiv1\pmod p\).  The geometry is canonical, but the Weyl phases are
not automatically compatible.

\subsection{Trace-Normalized Embeddings L3--T1}

Choose \(c_x=(c_y)_{y\mapsto x}\in B_x\) with
\(\Tr_{B_x/K}(c_x)=1\).
Such an element exists: for a finite separable field extension the trace
pairing is nondegenerate, so the trace map is a nonzero \(K\)-linear map to
the one-dimensional \(K\)-space \(K\), hence is surjective; products are
handled componentwise.

Define
\[
  \tau_{s,r,x}(q,p_0)=((q_y,c_y p_{0,y}))_{y\mapsto x}.
\]
For \(e=(q,p_0)\) and \(f=(q',p'_0)\), the fine exponent becomes
\[
\begin{aligned}
  \sum_{y\mapsto x}
  \Tr_{\kappa(y)/\mathbb F_p}
  \bigl(c_y(p_{0,y}q'_y-q_y p'_{0,y})\bigr)
  &=
  \Tr_{K/\mathbb F_p}
  \bigl(\Tr_{B_x/K}(c_x(a\otimes1))\bigr)\\
  &=
  \Tr_{K/\mathbb F_p}(a).
\end{aligned}
\]
Thus \(\tau_{s,r,x}\) preserves the character-valued Weyl commutator.  It is
injective because \(c_x\ne0\), so some component \(c_y\) is nonzero.

Taking the direct sum of \(\tau_{s,r,x}\) over the finite support of a label
gives
\[
  \tau_{s,r,U}:E_r(U)\to E_s(\pi_{s,r}^{-1}U).
\]
The Weyl relations are preserved on every finite support, so
\[
  \alpha_{s,r,U}(W_r(e))=W_s(\tau_{s,r,U}e)
\]
extends to a unital \(*\)-monomorphism
\[
  \alpha_{s,r,U}:\mathcal A_r(U)\hookrightarrow
  \mathcal A_s(\pi_{s,r}^{-1}U).
\]
For \(U\subset V\), these maps commute with the open-local inclusions because
both constructions are defined pointwise on closed supports.

\subsection{Canonical and Noncanonical Continuum Systems T2}

If \(p\nmid n=s/r\), then
\[
  c_x=n^{-1}\cdot1_{B_x}
\]
is canonical and has trace \(1\), because
\(\Tr_{B_x/K}(1_{B_x})=n\).  These choices are coherent: for
\(r\mid s\mid t\), the normalizing factors multiply as
\[
  (t/r)^{-1}=(s/r)^{-1}(t/s)^{-1}.
\]
Hence the scales \(r\) with \(p\nmid r\), ordered by divisibility, form a
canonical directed subsystem of Weyl nets.  Its algebraic continuum net is the
filtered colimit of these nets.

For the full tower of all \(r\), the current trace-character convention
requires extra data.  If \(p\mid n\), then \(n^{-1}\) does not exist.  A
trace-one element still exists, but it is not canonical in the above data.
The report therefore does not call the full all-rank tower canonical until a
coherent trace-one system, or an alternative compatible-character convention,
has been fixed.

\subsection{The Point over \texorpdfstring{\(\mathbb F_3\)}{F3} T3}

Let
\[
  X_0=\Spec\mathbb F_3,\qquad k_r=\mathbb F_{3^r}.
\]
Then \(X_r=\Spec k_r\) has one closed point and
\[
  E_r=k_r^2,\qquad
  \mathcal H_r=\ell^2(k_r),\qquad
  \mathcal A_r\cong M_{3^r}(\mathbb C).
\]
The character is
\[
  \psi_r(a)=\exp\!\left(\frac{2\pi i}{3}
  \Tr_{k_r/\mathbb F_3}(a)\right).
\]

For \(r\mid s\), put \(n=s/r\).  The naive inclusion
\((q,p)\mapsto(q,p)\) gives
\[
  \psi_s(pq'-q p')=\psi_r(pq'-q p')^n.
\]
Already \(r=1,s=2\) fails: \(\Tr_{\mathbb F_9/\mathbb F_3}(1)=2\), so the
phase of \(1\) becomes \(\exp(4\pi i/3)\), not \(\exp(2\pi i/3)\).
The canonical corrected map is
\[
  S_{s,r}(q,p)=(q,n^{-1}p)
\]
whenever \(3\nmid n\).  Consequently the canonical prime-to-\(3\) continuum
algebra of the arithmetic point is
\[
  \mathcal A_{\infty}^{(3')}
  =
  \varinjlim_{\substack{r\ge1\\3\nmid r}} M_{3^r}(\mathbb C),
\]
with transition maps \(W_r(q,p)\mapsto W_s(q,n^{-1}p)\).  This is an
algebraic colimit: every element is represented at a finite extension rank.
No Hilbert space \(\ell^2(\bar{\mathbb F}_3)\) or C*-completion is part of the
definition yet.

For the first missing full-tower step, \(r=1,s=3\), the naive trace is zero on
\(\mathbb F_3\) because \(\Tr_{\mathbb F_{27}/\mathbb F_3}(1)=3=0\).  A
trace-one element exists but is a choice.  For example, if
\(\theta^3+2\theta^2+1=0\), then this irreducible cubic has root trace
\(\Tr(\theta)=1\), and \((q,p)\mapsto(q,\theta p)\) is a valid embedding
\(\mathcal A_1\hookrightarrow\mathcal A_3\).  Choosing \(\theta\) is extra
structure, not a canonical consequence of \(k\subset k_3\).

\subsection{Frobenius on the Point Net T4}

At scale \(r\), define
\[
  F_r:E_r\to E_r,\qquad F_r(q,p)=(q^3,p^3).
\]
The trace is Frobenius-invariant:
\[
  \Tr_{k_r/\mathbb F_3}(a^3)
  =
  \sum_{i=0}^{r-1}a^{3^{i+1}}
  =
  \sum_{i=0}^{r-1}a^{3^i}
  =
  \Tr_{k_r/\mathbb F_3}(a).
\]
Therefore \(F_r\) preserves the Weyl commutator and induces
\[
  \Phi_r(W_r(q,p))=W_r(q^3,p^3).
\]
This is a Clifford/Gaussian automorphism of order \(r\).  For
\(k_2=\mathbb F_3[u]/(u^2+1)\), \(u^3=-u\), so
\[
  \Phi_2(W(a+bu,c+du))=W(a-bu,c-du).
\]

The prime-to-\(3\) embeddings are Frobenius-equivariant.  Indeed
\[
  S_{s,r}F_r(q,p)=(q^3,n^{-1}p^3)
  =
  F_sS_{s,r}(q,p),
\]
because \(n^{-1}\in\mathbb F_3\) is fixed by Frobenius.  Thus Frobenius acts
on the canonical prime-to-\(3\) continuum algebra
\(\mathcal A_{\infty}^{(3')}\).  In contrast, a noncanonical trace-one choice
for a \(3\)-divisible step is generally moved by Frobenius, so full-tower
Frobenius equivariance is additional structure.

\subsection{Conclusion}

Finite-field extension rank gives a precise arithmetic scale system.  Base
change of schemes is canonical, and closed points split by the finite
etale algebra \(\kappa(x)\otimes_{k_r}k_s\).  The Weyl net sees an extra
phase constraint: absolute trace characters force trace-normalized embeddings.
Therefore canonical continuum limits exist immediately on the prime-to-\(p\)
subtower; the full tower exists only after coherent trace-one or compatible
character data are added.  For the arithmetic point over \(\mathbb F_3\), the
canonical continuum object currently justified by the report is the
prime-to-\(3\) algebraic colimit of the matrix algebras \(M_{3^r}(\mathbb C)\),
and Frobenius acts on it by the compatible Clifford automorphisms
\(W(q,p)\mapsto W(q^3,p^3)\).
